\chapter{Qualifying-Examination Problems and Selected Solutions}
\label{ch:qualifying-problems}
\label{app:qualifying}

\section{Problem 1: substitution and satisfaction}
\label{sec:qualifying-problems-01}
\label{prob:substitution}

Let \(\varphi\) be a first-order formula, let \(x\) be a variable, let
\(t\) be a term, let
\(\Meta\) be a structure, and let \(s\) be a variable assignment.  Prove
\[
  \Meta,s\entails\varphi[t/x]
  \quad\Longleftrightarrow\quad
  \Meta,s[x\mapsto\sem{t}^{\Meta}_s]\entails\varphi.
\]

\begin{solution}
First prove the term-substitution lemma by structural induction on terms.  Then
proceed by structural induction on \(\varphi\).  Atomic formulas follow from
the term lemma.  For each connective, unfold its semantic clause and apply the
induction hypotheses to its immediate subformulas.  Let
\(Q\in\set{\forall,\exists}\) and consider \(Qy\,\psi\).  If \(y=x\),
substitution stops at the binder.  On the right side, the quantifier update of
\(x\) overrides the preceding update of \(x\), so the truth value of
\(Qx\,\psi\) is independent of that preceding value.

Now suppose \(y\neq x\).  If \(y\) occurs in \(t\), the deterministic
capture-avoiding definition first renames the binder to a variable fresh for
\(t\) and \(\psi\).  Thus it remains to consider
\(y\notin\operatorname{Var}(t)\).  For each \(a\in\abs{\Meta}\), the assignments
\[
  s[y\mapsto a][x\mapsto\sem{t}^{\Meta}_s]
  \quad\text{and}\quad
  s[x\mapsto\sem{t}^{\Meta}_s][y\mapsto a]
\]
agree because \(x\neq y\), and term evaluation of \(t\) is unchanged by the
update at \(y\).  Apply the induction hypothesis to \(\psi\) under each updated
assignment and then apply the semantic clause for \(Q\).  Alpha-invariance
justifies the preliminary binder renaming.  Substitution is total and
capture-avoiding, so the displayed definition requires no additional
definedness hypothesis.
\end{solution}

\section{Problem 2: compactness}
\label{sec:qualifying-problems-02}
\label{prob:compactness}

Let \(\Lang\) contain a binary relation symbol \(<\) and constants
\(c_n\) for \(n\in\N\).  Let \(T\) say that \(<\) is a strict linear order and
include \(c_{n+1}<c_n\) for every \(n\).  Prove that \(T\) has a model, and
explain why this does not prove that every strict linear order has an infinite
descending chain.

\begin{solution}
Every finite subset \(T_0\subseteq T\) mentions only finitely many of the
descending-chain axioms.  Interpret the relevant constants as a finite
descending segment of \(\Z\) with its usual order, and interpret unused
constants arbitrarily.  Then \(T_0\) has a model.  Compactness gives a model of
all of \(T\), in which the interpretations of the \(c_n\) form an infinite
descending chain.  The conclusion is existential: at least one strict linear
order with such a chain exists.  It makes no universal assertion about all
strict linear orders; \(\N\) with its usual order is a counterexample to that
universal assertion.
\end{solution}

\section{Problem 3: incompleteness hypotheses}
\label{sec:qualifying-problems-03}
\label{prob:incompleteness}

State a version of the first incompleteness theorem for a computably enumerable
theory \(T\), and explain why each hypothesis is needed.

\begin{solution}
One standard version is: if \(T\) is a consistent, computably enumerable
theory in the language of arithmetic extending a fixed weak theory sufficient
to represent the required computable syntactic relations, then there is a
sentence \(G_T\) such that \(T\notderives G_T\); with the stronger
\(\omega\)-consistency hypothesis in G\"odel's original form,
\(T\notderives\neg G_T\) as well.  Rosser's refinement replaces
\(\omega\)-consistency by consistency for a differently constructed sentence.
The computably enumerable axiom set makes provability arithmetically expressible;
arithmetic strength supports coding and diagonalization; consistency prevents
the theory from deriving every sentence; and the strengthened hypothesis in
the original proof excludes derivation of the negation.  No version licenses a
claim about an arbitrary collection of axioms lacking these properties.
\end{solution}

\section{Problem 4: modal frame condition}
\label{sec:qualifying-problems-04}
\label{prob:modal-frame}

Prove that \(\Box P\to P\) is valid on a frame iff the accessibility relation is
reflexive.

\begin{solution}
If the relation is reflexive and
\(\mathcal M,w\entails\Box P\), then every accessible world satisfies \(P\);
since \(w\) is accessible from itself, \(\mathcal M,w\entails P\).  Conversely,
suppose the frame is not reflexive at \(w\).  Let \(\mathcal M\) be the
structure on that frame whose valuation makes \(P\) true at every world
accessible from \(w\) and false at \(w\).  Then
\(\mathcal M,w\entails\Box P\) but \(\mathcal M,w\notentails P\), so the
formula is not valid on the frame.
\end{solution}

\section{Problem 5: cut and subformula property}
\label{sec:qualifying-problems-05}
\label{prob:cut-subformula}

Explain why cut elimination implies a subformula property for the cut-free
classical propositional sequent calculus, and state one consequence.

\begin{solution}
Inspect every propositional rule except cut.  Each formula in a premise is either copied from
the conclusion or is an immediate subformula of the principal formula in the
conclusion.  Induction upward through a cut-free derivation therefore shows
that every formula occurrence is a subformula of a formula in the end sequent,
apart from the declared treatment of structural copies.  Cut is the rule that
can introduce an arbitrary intermediate formula and then remove it.  A
consequence is consistency: a cut-free derivation of the empty sequent
would have a last rule, but no rule of the cut-free calculus---initial,
structural, or logical---has the empty sequent as its conclusion, since
every weakening or contraction conclusion contains a displayed formula
occurrence.
For the first-order calculus, quantifier premises contain substitution
instances.  Its analytic property must therefore be stated modulo term
substitution and eigenvariable renaming rather than as the literal subformula
property just proved.
\end{solution}

